
de façon que le
premier facteur permute \(T_1,\ldots,T_n\) et le second \(T_{n+1},\ldots,T_{2n}\)~; \(K[X_1,\ldots,X_{2n}]\) est plongé dans l'anneau des invariants en envoyant \(X_j\) sur la \(j\)-ième fonction symétrique des \(T_i\)~; \(f'\) envoie la \(i\)-ième fonction symétrique de \(T_1,\ldots,T_n\) (resp. \(T_{n+1},\ldots,T_{2n}\)) sur le \(i\)-ième coefficient de \(\chi_1\) (resp. \(\chi_2\)). Alors~:
\[
\begin{aligned}
\Phi(\chi)
&=H(f_\chi)\left(\Phi\left(\prod_{i=1}^{2n}(1+T_it)\right)\right)
=H(f'_\chi)\left(\Phi\left(\prod_{i=1}^{2n}(1+T_it)\right)\right);\\
\Phi(\chi_1) &=H(f'_\chi)\left(\Phi\left(\prod_{i=1}^{n}(1+T_it)\right)\right)\quad(\text{resp. }\chi_2,\ i=n+1,\ldots,2n).
\end{aligned}
\]
Comme \(H(f'_\chi)\) est un homomorphisme de groupes, on est ramené au cas où \(\chi=\prod_{i=1}^{2n}(1+T_it)\), et dans ce cas l'additivité de \(\Phi\) résulte de la définition des \(\Phi(1+X_1t+\ldots+X_k t^k)\).

1.7. On peut généraliser les considérations précédentes au cas d'un produit de foncteurs \(\Ghat\)~; le foncteur \((\Ghat)^k\) est représentable par l'anneau
\(K[(X_i^{(j)})]\), \(j=1,\ldots,k\), et \(i\geq 1\), où \(X_i^{(j)}\) est de poids \(i\). On notera \(\varphi^{(j)}\) la série \(1+X_1^{(j)}t+\ldots+X_n^{(j)}t^n+\ldots\)~; l'élément \((\varphi^{(j)})_j\) de \(\bigl(\Ghat(K[(X_i^{(j)})])\bigr)^k\) possède la propriété universelle analogue à celle de \(\varphi\) (cf. 1.2).

On peut introduire les morphismes \(f_n^{(j)}\), \(u_H^{n,(j)}\) comme précédemment en refaisant la construction pour chaque famille \((X_i^{(j)})_{i\geq 1}\) séparément. On obtient un résultat analogue au théorème 1.6 en cherchant les homomorphismes de foncteurs de \((\Ghat)^k\) dans \(H\), \(\mathbb Z\)-multilinéaires et compatibles aux filtrations par rapport à chacun des \(k\) arguments séparément : ils sont déterminés par leur valeur pour l'élément
